% Grabbed from https://tex.stackexchange.com/questions/85904/showcase-of-beautiful-title-page-done-in-tex

\renewcommand\epigraphflush{flushright}
\renewcommand\epigraphsize{\normalsize}
\setlength\epigraphwidth{0.7\textwidth}

\definecolor{titlepagecolor}{cmyk}{1,.60,0,.40}

\DeclareFixedFont{\titlefont}{T1}{ppl}{b}{it}{0.5in}

\makeatletter                       
\def\printauthor{%        
    {\large \@author}}              
\makeatother
\author{%
    Hallvard Enger Bjørgen\\
    Student \\
    \text{hallvard.bjorgen@gmail.com}\vspace{20pt} \\
    Amund Skavhaug \\
    Veileder \\
    \texttt{}
    }

% The following code is borrowed from: https://tex.stackexchange.com/a/86310/10898

\newcommand\titlepagedecoration{%
\begin{tikzpicture}[remember picture,overlay,shorten >= -10pt]

\coordinate (aux1) at ([yshift=-15pt]current page.north east);
\coordinate (aux2) at ([yshift=-410pt]current page.north east);
\coordinate (aux3) at ([xshift=-4.5cm]current page.north east);
\coordinate (aux4) at ([yshift=-150pt]current page.north east);

\begin{scope}[titlepagecolor!40,line width=12pt,rounded corners=12pt]
\draw
  (aux1) -- coordinate (a)
  ++(225:5) --
  ++(-45:5.1) coordinate (b);
\draw[shorten <= -10pt]
  (aux3) --
  (a) --
  (aux1);
\draw[opacity=0.6,titlepagecolor,shorten <= -10pt]
  (b) --
  ++(225:2.2) --
  ++(-45:2.2);
\end{scope}
\draw[titlepagecolor,line width=8pt,rounded corners=8pt,shorten <= -10pt]
  (aux4) --
  ++(225:0.8) --
  ++(-45:0.8);
\begin{scope}[titlepagecolor!70,line width=6pt,rounded corners=8pt]
\draw[shorten <= -10pt]
  (aux2) --
  ++(225:3) coordinate[pos=0.45] (c) --
  ++(-45:3.1);
\draw
  (aux2) --
  (c) --
  ++(135:2.5) --
  ++(45:2.5) --
  ++(-45:2.5) coordinate[pos=0.3] (d);   
\draw 
  (d) -- +(45:1);
\end{scope}
\end{tikzpicture}%
}

%% Title-page alternative 1:
% \begin{titlepage}
% \noindent
% \titlefont Forside 1: \\A prestudy of\\ intelligent edge video \\analytics devices;\\ best practices.\\ \par
% \epigraph{Division of images into several zones and only processing the parts of the images that are interesting could save processing time and cost and thus facilitate for more mobile edge-computing devices. This prestudy aims to challenge some of the limitations of pattern recognition tasks by exploring the use of predefined sectors in areas to push towards extended usability of edge-computing mobile devices.\vspace{0.5cm} 

% \mainRQ }%
% {\textit{Trondheim/Esbjerg 2023}\\ \textsc{Norwegian University of Science and Technology}}
% \null\vfill
% \vspace*{1cm}
% \noindent
% \hfill
% \begin{minipage}{0.35\linewidth}
%     \begin{flushright}
%         \printauthor
%     \end{flushright}
% \end{minipage}
% %
% \begin{minipage}{0.02\linewidth}
%     \rule{1pt}{125pt}
% \end{minipage}
% \titlepagedecoration
% \end{titlepage}


% Inspired by title template from ShareLaTeX Learn; Gubert Farnsworth & John Doe
% Edited by Jon Arnt Kårstad, NTNU IMT

\begin{titlepage}
\vbox{ }
\vbox{ }
\begin{center}
% Upper part of the page
\includegraphics[width=0.40\textwidth]{Images/NTNU_logo.png}\\[1cm]
\textsc{\LARGE Department of Mechanical and  }\\[0.3cm]
\textsc{\LARGE Industrial Engineering}\\[1.5cm]
\textsc{\Large Master's Thesis - robotics \& automation}\\[0.5cm]
\vbox{ }

% Title
% \HRule \\[0.4cm]
{ \huge \bfseries Efficient, accurate, and privacy-preservant object detection in edge devices}
\\[2.5cm] 
% Indoor vision transformer object detection accuracy augmentation by dataset specialization
% \\[0.5cm]

% { \LARGE \textit A prestudy of human localization on visual edge devices}
% \HRule \\
% Author
\large

% \begin{minipage}{0.35\linewidth}
%     \begin{flushright}
%         \printauthor
%     \end{flushright}
% \end{minipage}

\emph{Student:}\\
Hallvard Enger Bjørgen\\
\vspace{1cm}
\emph{Supervisor at NTNU:}\\
Amund Skavhaug
\vfill

% Bottom of the page
{\large Trondheim/Esbjerg Spring 2024}
\end{center}
\end{titlepage}